% Paper IV addendum: long-horizon assurance decomposition.
% Intended for insertion after the verified-checkpointing proposition.

\subsection{A long-horizon assurance decomposition}

The difference between an unchecked derivation and a verifier-gated research process can be stated as a simple event bound. Let $F$ be the event that a promoted result is false \emph{as the intended mathematical claim}. Let $E_{\mathrm{spec}}$ denote failure of the informal--formal specification alignment gate, and let $E_{\mathrm{trust}}$ denote failure of the proof-checking trust chain, including unsound checker behaviour, unregistered axioms, incorrect dependency identity or artifact substitution. Assume canonical promotion is fail-closed: a result is promoted only when the alignment witness and the exact proof receipt pass their respective gates.

\begin{proposition}[Assurance decomposition for a verifier-gated proof DAG]
Under the fail-closed promotion rule above,
\[
F\subseteq E_{\mathrm{spec}}\cup E_{\mathrm{trust}}.
\]
Consequently, for any probability measure describing the deployment environment,
\[
\Pr(F)\leq
\Pr(E_{\mathrm{spec}})+\Pr(E_{\mathrm{trust}}).
\]
In particular, local proposal errors made by the LLM do not appear as an additional additive term merely because the research horizon contains many generated steps, provided rejected proposals never enter the trusted proof DAG.
\end{proposition}

\begin{proof}
Consider a promoted result outside $E_{\mathrm{spec}}\cup E_{\mathrm{trust}}$. Since $E_{\mathrm{spec}}$ does not occur, the checked formal statement faithfully represents the intended claim under the declared alignment criterion. Since $E_{\mathrm{trust}}$ does not occur, the exact admitted proof artifact is valid under the registered assumptions and verifier semantics. Therefore the intended claim is not false, so $F$ cannot occur. This proves the set inclusion. The probability inequality follows from the union bound. \qed
\end{proof}

\begin{remark}
This is an assurance decomposition, not a calibrated numerical guarantee. Estimating either term is a separate empirical problem. The important structural point is that, once every accepted dependency is verifier-gated, a thousand hallucinated candidate moves need not create a thousand hidden logical liabilities. They can create substantial search cost, and they can indirectly increase pressure on formalization or infrastructure, but they acquire truth authority only by passing the explicit gates.
\end{remark}

For novelty, a distinct error event $E_{\mathrm{novel}}$ must be tracked. A result promoted to the strongest new-mathematics candidate stage therefore has a broader research-level risk decomposition,
\[
\begin{aligned}
\Pr(F_{\mathrm{research}})\leq{}&\Pr(E_{\mathrm{spec}})
+\Pr(E_{\mathrm{trust}})\\
&+\Pr(E_{\mathrm{novel}})
+\Pr(E_{\mathrm{value}}),
\end{aligned}
\]
where $E_{\mathrm{value}}$ denotes failure of the declared research-value review criterion. This second expression is likewise a union-bound decomposition, not an assumption that the events are independent.

\subsection{Obstruction--transformation routing is a discovery-process gate, not theorem authority}

The deployed mathematical-research workflow now inserts an additional fail-closed proposal layer before candidate generation. Once an atomic obstruction and its context fibre are frozen, the strict path is
\[
\begin{aligned}
&\text{context/method-transfer review}
\rightarrow \text{success+failure experience review}\\
&\rightarrow \text{content-bound obstruction--transformation review}
\rightarrow \text{hash-chained next-step trace}\\
&\rightarrow \text{candidate proposal and falsification}.
\end{aligned}
\]
The transformation memory stores scoped source episodes $O\xrightarrow{T}O'$ with explicit source preconditions, observed effects, preserved invariants, broken constraints, breakpoints, provenance and source authority. The active target is fingerprinted by roles, relations, constraints, failure mechanisms, invariants to preserve, desired transition and forbidden losses rather than by topic vocabulary alone.

The router follows the invention-last order
\[
\texttt{SEARCH}\rightarrow\texttt{JUMP}\rightarrow\texttt{GLUE}\rightarrow\texttt{LIFT},
\]
with \texttt{CANNOT\_CHECK} when memory identity, mapping, composition, coverage or residual evidence is incomplete. Direct and cross-domain single-episode routes require complete desired-effect coverage. A cross-domain source additionally needs an explicit structural/precondition mapping witness. Verified partial effects are reserved for \texttt{GLUE}, whose composition witness must bind operation order, interfaces, incompatibility checks and target validation.

\texttt{LIFT} is especially constrained because a failed proof attempt is weak evidence for a missing mathematical primitive. It is available only after SEARCH/JUMP/GLUE are explicitly bounded, retrieved candidates are accounted for, a cross-problem coverage receipt supports the no-match claim, and at least two distinct failed attempts share residual structure. The output is a \emph{MissingTransformationSpecification} describing what a later representation or operator must preserve, break, expose, reduce and validate. It neither asserts that such an object exists nor promotes the object if one is proposed.

These gates affect proposal provenance and search discipline, not the assurance decomposition above. In particular,
\[
\text{valid source transformation}
\not\Rightarrow
\text{valid target transfer}
\not\Rightarrow
\text{theorem truth}
\not\Rightarrow
\text{novelty}.
\]
A candidate produced by SEARCH, JUMP, GLUE or LIFT still proceeds through the paper's ordinary specification-alignment, falsification, formal proof, verifier-trust, bounded-novelty and research-value gates. Conversely, a machine-checked theorem can be true even if the strict RAKL discovery chronology was not followed; theorem assurance does not retroactively certify discovery-process compliance.

The semantic-shortcut implementation therefore extends the paper's discovery architecture without changing what a proof receipt means. Its empirical effect on proof/search efficiency or success is a separate matched-evaluation question registered outside the theorem-assurance authority chain.
